% REFERÃÅ NCIAS------------------------------------------------------------------
\usepackage[%
    alf,
    abnt-emphasize=bf,
    bibjustif,
    recuo=0cm,
    abnt-url-package=url,       % Utiliza o pacote url
    abnt-refinfo=yes,           % Utiliza o estilo bibliográfico abnt-refinfo
    abnt-etal-cite=3,
    abnt-etal-list=3,
    abnt-thesis-year=final
]{abntex2cite}                  % Configura as citações bibliográficas conforme a norma ABNT

% PACOTES----------------------------------------------------------------------
\usepackage[utf8]{inputenc}                                 % Codificação do documento
\usepackage[T1]{fontenc}                                    % Seleção de código de fonte
\usepackage{booktabs}                                       % Réguas horizontais em tabelas
\usepackage{color, colortbl}                                % Controle das cores
\usepackage{float}                                          % Necessário para tabelas/figuras em ambiente multi-colunas
\usepackage{graphicx}                                       % Inclusão de gráficos e figuras
\usepackage{indentfirst}                                    % Indenta o primeiro parágrafo de cada seção
\usepackage[protrusion=true,expansion=false]{microtype}     % Melhora a justificação do texto (ABNT)
\usepackage{array}                                          % Especificadores de coluna avançados (>{...}) em tabelas
\usepackage{tabularx}                                       % Tabelas com colunas de largura automática (preenchem \textwidth)
\usepackage{lastpage}                                       % Para encontrar última página do documento
\usepackage{amsfonts, amssymb, amsmath}                     % Fontes e símbolos matemáticos
\usepackage[scaled]{helvet}                                 % Usa a fonte Helvetica (equivalente Arial)
%\usepackage{times}                                         % Usa a fonte Times
%\usepackage{palatino}                                      % Usa a fonte Palatino
%\usepackage{lmodern}                                       % Usa a fonte Latin Modern
\usepackage[bottom]{footmisc}                               % Mantém as notas de rodapé sempre na mesma posição
\usepackage{latexsym}                                       % Símbolos matemáticos
\usepackage{xurl}                                           % Quebra URLs longas em qualquer ponto (evita estouro de margem nas referências)
%\usepackage{picinpar}                                      % Dispor imagens em parágrafos
%\usepackage{scalefnt}                                      % Permite redimensionar tamanho da fonte
%\usepackage{subfig}                                        % Posicionamento de figuras
%\usepackage{upgreek}                                       % Fonte letras gregas''
\usepackage{ifthen}
\usepackage[pass]{geometry}                                 % Permite \newgeometry/\restoregeometry sem alterar o layout do memoir
\usepackage{tikz}
\usetikzlibrary{positioning, arrows.meta, shapes.geometric, calc, decorations.pathreplacing}

% Tipo de float "Gráfico" (sintaxe do pacote float, já carregado)
\newcommand{\listofgraficosname}{LISTA DE GRÁFICOS}
\newfloat{grafico}{htbp}{lgr}[chapter]
\floatname{grafico}{Gráfico}


% Redefine a fonte para uma fonte similar a Arial (fonte Helvetica)
\renewcommand*\familydefault{\sfdefault}

% Impede páginas em branco entre capítulos e elementos pré-textuais
\let\cleardoublepage\clearpage
\let\clearforchapter\clearpage

% ABNT NBR 6024 — FORMATAÇÃO DE TÍTULOS DE SEÇÃO
% Primária   (\chapter):       CAIXA ALTA + negrito
% Secundária (\section):       CAIXA ALTA + sem negrito
% Terciária  (\subsection):    caixa baixa + negrito
% Quaternária(\subsubsection): caixa baixa + sem negrito
% Quinária   (\paragraph):     caixa baixa + itálico
\setboolean{ABNTEXupperchapter}{true}
\setboolean{ABNTEXuppersection}{true}
\newcommand{\ABNTsecformat}[1]{\normalfont\MakeUppercase{#1}}
\setsecheadstyle{\ABNTsecformat}
\setsubsecheadstyle{\normalfont\bfseries}
\setsubsubsecheadstyle{\normalfont}
\setparaheadstyle{\normalfont\itshape}
% Sumário: entradas de capítulo e seção em caixa alta
\makeatletter
\AtBeginDocument{%
  \let\ABNTEXorig@l@chapter\l@chapter
  \renewcommand{\l@chapter}[2]{\ABNTEXorig@l@chapter{\MakeUppercase{#1}}{#2}}%
  \let\ABNTEXorig@l@section\l@section
  \renewcommand{\l@section}[2]{\ABNTEXorig@l@section{\MakeUppercase{#1}}{#2}}%
}
\makeatother

% Sumário SEM RECUO (replica a opção "semrecuonosumario" do modelo oficial UTFPR):
% todos os níveis alinhados à margem esquerda, conforme NBR 6027:2012.
\setlength{\cftchapterindent}{0pt}
\setlength{\cftsectionindent}{0pt}
\setlength{\cftsubsectionindent}{0pt}
\setlength{\cftsubsubsectionindent}{0pt}
\setlength{\cftparagraphindent}{0pt}

% Sumário — estilo por nível espelhando a NBR 6024 (mesma formatação do texto):
%   nível 1 (capítulo):     CAIXA ALTA + negrito  (classe + \l@chapter)
%   nível 2 (seção):        CAIXA ALTA + sem negrito (\l@section faz a caixa alta)
%   nível 3 (subseção):     caixa baixa + negrito
%   nível 4 (subsubseção):  caixa baixa + sem negrito
\renewcommand{\cftsectionfont}{\normalfont}
\renewcommand{\cftsubsectionfont}{\normalfont\bfseries}
\renewcommand{\cftsubsubsectionfont}{\normalfont}


% Listagens de código (usadas no apêndice de consultas SQL). `listings` é
% suficiente aqui: as listagens são curtas e não precisam de realce semântico
% completo — `minted` exigiria Pygments e compilação com -shell-escape.
\usepackage{listings}
\lstset{
    basicstyle=\ttfamily\scriptsize,
    keywordstyle=\bfseries,
    commentstyle=\itshape\color{black!55},
    stringstyle=\color{black!70},
    numbers=none,
    breaklines=true,
    breakatwhitespace=true,
    showstringspaces=false,
    frame=single,
    framesep=4pt,
    rulecolor=\color{black!40},
    captionpos=b,
    % Acentos dentro das listagens: sem isto, caracteres UTF-8 saem quebrados.
    extendedchars=true,
    inputencoding=utf8,
    literate={á}{{\'a}}1 {â}{{\^a}}1 {ã}{{\~a}}1 {à}{{\`a}}1
             {é}{{\'e}}1 {ê}{{\^e}}1 {í}{{\'i}}1 {ó}{{\'o}}1
             {ô}{{\^o}}1 {õ}{{\~o}}1 {ú}{{\'u}}1 {ç}{{\c c}}1
}
\renewcommand{\lstlistingname}{Código}
